\par\noindent\textbf{(1)} Hamilton量为 \(H=p_r^2/(2m)+L^2/(2mr^2)-k/r\)。令 \(S=W_r(r)+L\phi-Et\)，得到 \(p_r=W_r'=\sqrt{2mE+2mk/r-L^2/r^2}\)。其两个正根界定束缚径向区间。
\par\noindent\textbf{(2)} 径向作用 \(J_r=(2\pi)^{-1}\oint p_r\,dr\)。对开普勒势积分可得 \(J_r=mk/\sqrt{-2mE}-L\)。每项量纲均为角动量；极圆轨道 \(J_r=0\)，此时 \(L=mk/\sqrt{-2mE}\)。
\par\noindent\textbf{(3)} 反解得到 \(H(J_r,L)=-mk^2/[2(J_r+L)^2]\)。因此 \(\omega_r=\partial H/\partial J_r=mk^2/(J_r+L)^3\)，\(\omega_\phi=\partial H/\partial L\) 同值。径向和角向频率恰相等，解释理想Kepler椭圆无拱点进动。
\par\noindent\textbf{(4)} 若增加小势 \(+\alpha/r^2\)，等效离心项由 \(L^2/(2m)\) 改为 \(L^2/(2m)+\alpha\)，可定义 \(L_{\mathrm{eff}}=\sqrt{L^2+2m\alpha}\)。作用变成 \(J_r=mk/\sqrt{-2mE}-L_{\mathrm{eff}}\)，故 \(\omega_\phi/\omega_r=L/L_{\mathrm{eff}}\) 的倒数需由实际作用偏导计算；正确为 \(\omega_\phi/\omega_r=L/L_{\mathrm{eff}}\)，表明拱点进动。
